% Section 4. Sources: 1367 PROOF_LOWER.md sections 2-4; logs/check_flow_{a,b}_r10.log, logs/float_table.log;
% Paper B sections 5-6 (Theorem F).
%%%%%%%%%%%%%%%%%%%%%%%%%%%%%%%%%%%%%%%%%%%%%%%%%%%%%%%%%%%%%%%%%%%%%%%%
\section{The lower bound}\label{sec:lower}
%%%%%%%%%%%%%%%%%%%%%%%%%%%%%%%%%%%%%%%%%%%%%%%%%%%%%%%%%%%%%%%%%%%%%%%%

\subsection{Stationary random rules}

For a lower bound we must exhibit many closed walks. We let the machine run on infinitely deep stacks and choose
each move at random by a rule that looks only at the top of the stacks.

\emph{Infinite stacks and windows.} Let $\Omega$ be the set of triples $\omega=(U,L,c)$ with $U=(U_1,U_2,\ldots)$ and
$L=(L_1,L_2,\ldots)$ infinite sequences over $\{0,1\}$, indexed from the top, and $c$ a parity. The four moves act
on $\Omega$ as before; a pop needs top symbol $1-c$. Fix a radius $r\ge1$. The \emph{window} of $\omega$ is
$w=(U_1\ldots U_r,\;L_1\ldots L_r,\;c)$.

\emph{The reference law.} Below the window the stacks are filled by a fixed background law $\beta$: a stationary
Markov chain on $\{0,1\}$, read downward, of memory $m\le r$, with all transition probabilities positive and
\emph{reversible}, that is, its law is unchanged when finite blocks are read backward. We use the simplest case: the
next symbol differs from the previous one with probability $s$ and equals it with probability $1-s$, with $0<s<1$
(a two-state chain is always reversible). For a word $w_1\ldots w_k$ with $k\ge m$, its \emph{continuation} is the
conditional law of the following symbols given $w_1\ldots w_k$.

\emph{Flows.} A \emph{flow} is a family of numbers $J(w,\xi)\ge0$ with total $1$, one for each window $w$ and each move
$\xi$ allowed at $w$. Put $\nu(w)=\sum_\xi J(w,\xi)$ and $p(\xi\mid w)=J(w,\xi)/\nu(w)$. The flow defines a random
infinite state: the window has law $\nu$, and the two stacks continue below it as independent continuations under
$\beta$. It also defines a random move, chosen with probability $p(\xi\mid w)$. Intuitively, $J(w,\xi)$ is the
probability of seeing window $w$ and then making move $\xi$. The conditions are:
\begin{itemize}
\item[(S)] \emph{Stationarity.} For every \emph{extended window} $Y$, with $r+1$ symbols of each stack, the
probability that the extended window after one move is $Y$ equals its probability now, which is $\nu(Y^\circ)$ times
the reference probabilities of the two symbols at depth $r+1$ given the symbols above them; $Y^\circ$ is $Y$ without
those two symbols.
\item[(D)] \emph{Zero drift.} The total flow of $\pA$ equals that of $\qA$, and the same for $L$.
\item[(P)] \emph{Positivity.} $J(w,\pA)>0$ and $J(w,\pB)>0$ for every window $w$.
\end{itemize}

\begin{theoremL}
If a flow satisfies (S), (D) and (P), then $\mu\ge e^{h}$, that is, $\mu^2\ge e^{2h}$, where
\[ h\;=\;\sum_{w,\xi}J(w,\xi)\,\log\frac{\nu(w)}{J(w,\xi)} \]
is the entropy of one move given the window.
\end{theoremL}

The proof follows that of Theorem~F of Paper~B~\cite{PaperB} in three steps.

\emph{Step 1: stationarity.} After a push the symbols below the new extended window are old symbols that continued
the old window; after a pop the two newly revealed symbols are fresh continuations. So, given the new extended
window, the rest is again a continuation (this uses $m\le r$), and by (S) the law of the new extended window is the
old one. Hence the random infinite state has the same law after every move.

\emph{Step 2: ergodicity.} Let $f$ be a bounded function with $\Ex[f(\omega_1)\mid\omega_0]=f(\omega_0)$. Stationarity gives
$\Ex(f(\omega_1)-f(\omega_0))^2=0$, so $f$ is unchanged by every move of positive probability, in particular by both
pushes, by (P). Comparing the two pushes from parity $1$ gives $f(1U,L,0)=f(U,1L,0)$.
For the other symbol, use the three comparisons
\[ f(0U,L,0)=f(0U,0L,1)=f(U,00L,1)=f(U,0L,0). \]
The first and last compare a push on $L$; the middle compares the two pushes from $(U,0L,0)$.
Thus $f(xU,L,0)=f(U,xL,0)$ for either symbol: moving a top symbol between stacks does not change $f$. Glue the two stacks into one two-sided sequence
$(\ldots,U_2,U_1\mid L_1,L_2,\ldots)$; then $f$ is invariant under shifting this sequence. Because $\beta$ is
reversible with positive transitions, the glued sequence has the same null sets as the stationary two-sided chain,
which is ergodic under the shift. So $f$ is constant almost surely, and the process is ergodic.

\emph{Step 3: counting.} By the ergodic theorem and (D), in $t$ moves both stack heights change by $o(t)$, and the
product of the probabilities of the chosen moves is about $e^{-ht}$. So from some fixed infinite state there are at
least $\frac12e^{t(h-\varepsilon)}$ move sequences of length $t$ along which no stack ever drops by more than
$\varepsilon t$. They only read the top $\varepsilon t+1$ symbols of each stack. Copy those symbols into a finite
state, adding below them a run of equal symbols so that the charge rule of Lemma~\ref{lem:words} holds. By
Lemma~\ref{lem:sweep} this state is reached from $e$, and every end state returns to $e$, in $O(\varepsilon t)$ moves.
This gives $e^{t(h-\varepsilon)}/O(t)$ distinct closed walks of one length between $t$ and $t+O(\varepsilon t)$, and
$z_n\le\mu^{2n}$ gives $h\le\log\mu$ as $t\to\infty$ and $\varepsilon\to0$. \qed

\subsection{Reversible flows}

For meanders, the hard part was to make (S) hold exactly: it is a large linear system, and Paper~B solved it at a
small radius and corrected the solution with about a million closed-walk directions. For the two-stack machine there
is a direct construction. Write $\beta(a)$ for the stationary probability of a word $a$ (read from the top),
$\beta(w)=\beta(U_1\ldots U_r)\,\beta(L_1\ldots L_r)/2$ for a window $w$, and $q(x\mid w)$ for the probability that
the symbol above the top of $U$ is $x$ given the window; for our chain $q(x\mid w)$ is $s$ if $x\neq U_1$ and $1-s$
otherwise. Let $w^U$ be the window after $\pA$ at $w$.

\begin{proposition}\label{prop:rev}
Let $X_U(w)>0$ be given for every window with $U_1\ne c$, and $X_L(w)>0$ for every window with $L_1\ne c$. Put
\[ J(w,\qA)=\beta(w)\,X_U(w),\qquad J(w,\pA)=\beta(w)\,q(c\mid w)\,X_U(w^U), \]
and the same for $L$. After normalization, $J$ satisfies (S), (D) and (P).
\end{proposition}

\begin{proof}
The random infinite state of the flow has density $\nu(w)/\beta(w)$ with respect to the product of two stationary
chains and a fair parity. The move $\pA$ maps $\omega$ to $\omega'$ with one more symbol, $c$, on top of $U$, and the
reference probability is multiplied by $q(c\mid w)$, because for a reversible chain the probability of the symbol
above equals that of the symbol below. The flow forward along this move is $J(w,\pA)/\beta(w)$ per unit of reference
probability at $\omega$, that is $q(c\mid w)X_U(w^U)$; the flow back along $\qA$ from $\omega'$ is
$J(w^U,\qA)/\beta(w^U)=X_U(w^U)$ per unit of reference probability at $\omega'$, and the reference probability at
$\omega'$ is $q(c\mid w)$ times that at $\omega$. So the two flows are equal. Every move is undone by a move, so at every infinite state the total inflow equals
the total outflow, and (S) follows by summing over the states with a given extended window. Summing the equality over
all pushes on $U$ gives (D). (P) holds since every $X>0$.
\end{proof}

Intuitively, the rule is in detailed balance: in the long run it crosses every move as often forward as backward.
Such a rule is automatically stationary and automatically has zero drift. Any positive numbers $X$ give a valid lower
bound; optimization only chooses them well. At an optimum, $\log X(v)$ is a weighted average of $\log(\nu/J)$ over
the flow values that use $X(v)$, and iterating this relation converges quickly.

\subsection{The certificate}

We take $s=41/50$: the reference stacks prefer to alternate, as typical stacks do. With $s=1/2$ (independent
symbols) the bounds are smaller, $9.983$ instead of $10.084$ at radius $9$. The numbers $X$ were scaled so that the
largest is $2^{60}$ and rounded to integers; then every flow value is an integer after a common scaling, and only the
entropy needs rounding. Table~\ref{tab:lower} gives the bound by radius.

\begin{table}[htb]
\centering
\begin{tabular}{rrrl}
\toprule
radius $r$ & windows & $e^{2h}$ & status\\
\midrule
$1$ & $8$ & $9.497$ & float\\
$2$ & $32$ & $9.880$ & float\\
$4$ & $512$ & $10.030980878326$ & certified\\
$6$ & $8{,}192$ & $10.065$ & float\\
$8$ & $131{,}072$ & $10.079276064125$ & certified\\
$10$ & $2{,}097{,}152$ & $\mathbf{10.086918840819}$ & certified\\
\bottomrule
\end{tabular}
\caption{Lower bounds $\mu^2\ge e^{2h}$ from reversible flows with $s=41/50$. Certified values are rounded down and
were accepted by two checkers and by a third written in an independent reading; float values are from the
optimizer and are not bounds.}
\label{tab:lower}
\end{table}

At radius $10$ the certificate has $2{,}097{,}152$ windows and $6{,}291{,}456$ flow values, and
\[ h\;\ge\;1.15561971006089,\qquad \mu^2\;\ge\;e^{2h}\;\ge\;10.086918840819, \]
the lower half of Theorem~\ref{thm:main}. Already radius $1$, a rule that sees only the top symbol of each stack,
gives about $9.497$ in floating point, well above $8$. The bound grows slowly with the radius, by $0.0076$ from radius $8$ to $10$, so the
gap of $0.026$ to the estimate would need much larger windows.
